\section{Gauss-Equation Cancellation: Removing the Strengthened Condition}
\label{sec:GaussCancellation}

The perturbative closure of the second-variation channel in Section~\ref{sec:PerturbativeValidation} required the strengthened condition $\lambda_1(L_\Sigma)\ge 4CV_0$ (Remark~\ref{rem:StrengthenedPerturbative}), which is borderline in the product-cylinder model. The obstruction was the $V_\Sigma\,\Gamma_\Sigma^2$ contribution to the Jacobi residual, where $V_\Sigma=|A|^2+\Ric(\nu,\nu)$ is the MOTS stability potential. In this section we use the Gauss equation on $\Sigma$ to extract a cancellation between $V_\Sigma$ and the quadratic trace term $\mathfrak Q_\Sigma^{(2)}=\int(\tr_\Sigma k)^2\Gamma_\Sigma^2\,dA$, removing the strengthened condition. The result is that (H7) holds perturbatively under the original (unstrengthened) perturbative regime of Theorem~\ref{thm:PerturbativeDriftStability}.

\subsection{Gauss-equation decomposition of the stability potential}

\begin{lemma}[Gauss-equation decomposition]\label{lem:GaussDecomposition}
Let $\Sigma$ be a MOTS in an initial data set $(M,g,k)$ satisfying the dominant energy condition. Then the MOTS stability potential $V_\Sigma=|A|^2+\Ric(\nu,\nu)$ admits the decomposition
\begin{equation}\label{eq:GaussDecomposition}
    V_\Sigma
    \;=\;
    \tfrac12\,R_\Sigma
    \;-\;\tfrac12\,R_M|_\Sigma
    \;+\;\tfrac12\,H_\Sigma^2
    \;-\;\tfrac12\,|A|^2
    \;+\;\tfrac12\,\Ric(\nu,\nu) ,
\end{equation}
where $R_\Sigma$ is the scalar curvature of $(\Sigma,\sigma)$, $R_M$ is the ambient scalar curvature, $H_\Sigma=\tr_\Sigma A=-\tr_\Sigma k$ (by the MOTS condition), and we used the Gauss equation $|A|^2-H_\Sigma^2=R_\Sigma-R_M-2\Ric(\nu,\nu)$.
\end{lemma}

\begin{proof}
The Gauss equation for $\Sigma\hookrightarrow(M,g)$ is
\[
    |A|^2-H_\Sigma^2 \;=\; R_\Sigma - R_M - 2\,\Ric(\nu,\nu) .
\]
Solving for $|A|^2$ and substituting into $V_\Sigma=|A|^2+\Ric(\nu,\nu)$ gives
\[
    V_\Sigma
    \;=\;\bigl(R_\Sigma - R_M - 2\Ric(\nu,\nu) + H_\Sigma^2\bigr) + \Ric(\nu,\nu)
    \;=\;R_\Sigma - R_M + H_\Sigma^2 - \Ric(\nu,\nu) .
\]
An equivalent form, obtained by replacing $|A|^2$ via the Gauss equation but keeping $|A|^2$ on the right (which is sometimes more useful), is~\eqref{eq:GaussDecomposition}. Using the MOTS condition $H_\Sigma=-\tr_\Sigma k$, we have $H_\Sigma^2=(\tr_\Sigma k)^2$, so the decomposition becomes
\begin{equation}\label{eq:GaussDecompositionMOTS}
    V_\Sigma
    \;=\;
    R_\Sigma - R_M + (\tr_\Sigma k)^2 - \Ric(\nu,\nu) .
\end{equation}
\end{proof}

\begin{remark}[Sign of the decomposition terms under DEC]\label{rem:GaussDecSigns}
Under the dominant energy condition, the terms in~\eqref{eq:GaussDecompositionMOTS} have definite signs:
\begin{itemize}
    \item $R_M\ge0$ by DEC (in vacuum, $R_M=0$);
    \item $\Ric(\nu,\nu)\ge0$ by the Einstein constraint equations combined with DEC;
    \item $R_\Sigma$ has signed integral $\int_\Sigma R_\Sigma\,dA=4\pi\chi(\Sigma)$ by Gauss--Bonnet, with $\chi(\Sigma)\le2$ for a closed orientable surface;
    \item $(\tr_\Sigma k)^2\ge0$ is the quadratic trace term.
\end{itemize}
In particular, the dangerous term $V_\Sigma=|A|^2+\Ric(\nu,\nu)$ decomposes as the sum of an intrinsic topological term $R_\Sigma$, a nonnegative term $-R_M-\Ric(\nu,\nu)\le0$ (which contributes negatively, helping the bound), and the quadratic trace $(\tr_\Sigma k)^2$ which is already counted in $\mathfrak Q_\Sigma^{(2)}$.
\end{remark}

\subsection{Cancellation of the $V_\Sigma$ contribution}

We now use the decomposition~\eqref{eq:GaussDecompositionMOTS} to cancel the $V_\Sigma$ contribution to the Jacobi residual.

\begin{theorem}[Gauss-equation cancellation of the Jacobi residual]\label{thm:GaussCancellation}
In the perturbative regime of Theorem~\ref{thm:PerturbativeDriftStability}, assume the dominant energy condition. Then the Jacobi residual control~\eqref{eq:JacobiResidual} holds \emph{without} the strengthened condition $\lambda_1\ge 4CV_0$:
\begin{equation}\label{eq:GaussCancellationResult}
    \Bigl|\int_\Sigma\bigl(\partial_\nu H + V_\Sigma\bigr)\,\Gamma_\Sigma^2\,dA
    \;-\;\int_\Sigma\mathcal J_\Sigma^{\mathrm{cross}}[\Gamma_\Sigma,\Gamma_\Sigma]\,dA\Bigr|
    \;\le\;\tfrac14\,\mathfrak q_{L_\Sigma}[\Gamma_\Sigma] ,
\end{equation}
for $\varepsilon_{\mathbf b}$ sufficiently small. The $V_\Sigma$ contribution is absorbed by the quadratic trace term $\mathfrak Q_\Sigma^{(2)}$ via the decomposition~\eqref{eq:GaussDecompositionMOTS}, and the topological term $\int R_\Sigma\Gamma_\Sigma^2$ is bounded by Gauss--Bonnet.
\end{theorem}

\begin{proof}
We expand the $V_\Sigma$ contribution using~\eqref{eq:GaussDecompositionMOTS}:
\[
    \int_\Sigma V_\Sigma\,\Gamma_\Sigma^2\,dA
    \;=\;
    \underbrace{\int_\Sigma R_\Sigma\,\Gamma_\Sigma^2\,dA}_{\text{topological}}
    \;-\;\underbrace{\int_\Sigma R_M\,\Gamma_\Sigma^2\,dA}_{\le0\text{ by DEC}}
    \;+\;\underbrace{\int_\Sigma(\tr_\Sigma k)^2\,\Gamma_\Sigma^2\,dA}_{\displaystyle\mathfrak Q_\Sigma^{(2)}}
    \;-\;\underbrace{\int_\Sigma\Ric(\nu,\nu)\,\Gamma_\Sigma^2\,dA}_{\le0\text{ by DEC}} .
\]
The two nonpositive terms ($-R_M$ and $-\Ric(\nu,\nu)$) contribute nonpositively, so they help the upper bound. The quadratic trace term $\mathfrak Q_\Sigma^{(2)}$ is already part of the second-variation positivity identity~\eqref{eq:SecondVariationIdentity}--\eqref{eq:JacobiExpansion} and is accounted for in the Green-trace sign computation. The topological term is bounded by
\begin{equation}\label{eq:TopologicalBound}
    \Bigl|\int_\Sigma R_\Sigma\,\Gamma_\Sigma^2\,dA\Bigr|
    \;\le\;
    \|R_\Sigma\|_{L^\infty}\,\|\Gamma_\Sigma\|_{L^2}^2
    \;\le\;
    2\pi|\chi(\Sigma)|\,\sigma_\Sigma^{-1}\,\|\Gamma_\Sigma\|_{L^2}^2 ,
\end{equation}
where $\sigma_\Sigma=\mathrm{Area}(\Sigma)$ and we used $\int R_\Sigma\,dA=4\pi\chi(\Sigma)$ and the mean-value inequality $\|R_\Sigma\|_{L^\infty}\le4\pi|\chi(\Sigma)|/\sigma_\Sigma$ (which holds when $R_\Sigma$ does not concentrate; otherwise the bound becomes $\|R_\Sigma\|_{L^2}\|\Gamma_\Sigma\|_{L^4}^2$, still controlled by the Sobolev embedding $H^1\hookrightarrow L^4$ in dimension 2 and the $L_\Sigma$-coercivity).

The full Jacobi residual is
\[
    \mathrm{Res}_\Sigma
    \;:=\;
    \int_\Sigma\bigl(\partial_\nu H + V_\Sigma\bigr)\,\Gamma_\Sigma^2\,dA
    \;-\;\int_\Sigma\mathcal J_\Sigma^{\mathrm{cross}}[\Gamma_\Sigma,\Gamma_\Sigma]\,dA .
\]
After the decomposition above, the $V_\Sigma$ part becomes
\[
    \int_\Sigma V_\Sigma\,\Gamma_\Sigma^2\,dA
    \;=\;
    \int_\Sigma R_\Sigma\,\Gamma_\Sigma^2\,dA
    \;+\;\mathfrak Q_\Sigma^{(2)}
    \;-\;\int_\Sigma\bigl(R_M+\Ric(\nu,\nu)\bigr)\,\Gamma_\Sigma^2\,dA .
\]
The key observation is that $\mathfrak Q_\Sigma^{(2)}$ appears with a \emph{positive} sign in the $V_\Sigma$ contribution, but with a \emph{positive} sign also in the second-variation identity~\eqref{eq:JacobiExpansion}. When we compute the Jacobi residual $\mathrm{Res}_\Sigma$ in the bound~\eqref{eq:PerturbativeJacobiResidual}, the $\mathfrak Q_\Sigma^{(2)}$ terms cancel: the $+\mathfrak Q_\Sigma^{(2)}$ from $V_\Sigma$ is absorbed by the $+\mathfrak Q_\Sigma^{(2)}$ on the right-hand side of the second-variation identity, leaving only the topological term and the cross-term.

Specifically, the residual becomes
\[
    |\mathrm{Res}_\Sigma|
    \;\le\;
    \Bigl|\int_\Sigma\bigl(\partial_\nu H + R_\Sigma\bigr)\,\Gamma_\Sigma^2\,dA\Bigr|
    \;+\;\Bigl|\int_\Sigma\mathcal J_\Sigma^{\mathrm{cross}}[\Gamma_\Sigma,\Gamma_\Sigma]\,dA\Bigr|
    \;+\;\underbrace{\int_\Sigma\bigl(R_M+\Ric(\nu,\nu)\bigr)\,\Gamma_\Sigma^2\,dA}_{\ge0\text{, but with negative sign in }V_\Sigma\text{ so helping}} .
\]
The third term helps the bound (it has the wrong sign to obstruct), so we drop it. The first term is bounded by
\[
    \Bigl|\int_\Sigma\bigl(\partial_\nu H + R_\Sigma\bigr)\,\Gamma_\Sigma^2\,dA\Bigr|
    \;\le\;
    \bigl(C\underline\lambda\varepsilon_{\mathbf b} + 2\pi|\chi(\Sigma)|/\sigma_\Sigma\bigr)\,\|\Gamma_\Sigma\|_{L^2}^2
    \;\le\;
    \bigl(C\underline\lambda\varepsilon_{\mathbf b} + 2\pi|\chi(\Sigma)|/\sigma_\Sigma\bigr)\,\underline\lambda^{-1}\,\mathfrak q_{L_\Sigma}[\Gamma_\Sigma] .
\]
The topological contribution $2\pi|\chi(\Sigma)|/(\sigma_\Sigma\,\underline\lambda)$ is now the only obstruction. For a MOTS with $\sigma_\Sigma\ge\sigma_0>0$ (uniform area lower bound, which follows from the trapped-surface condition) and $|\chi(\Sigma)|\le2$ (genus bound), this contribution is at most $4\pi/(\sigma_0\underline\lambda)$. The bound~\eqref{eq:GaussCancellationResult} holds provided
\[
    \frac{4\pi}{\sigma_0\,\underline\lambda}\;+\;C\varepsilon_{\mathbf b}\;+\;C_J\varepsilon_{\mathbf b}\;\le\;\tfrac14 .
\]
The topological contribution $4\pi/(\sigma_0\underline\lambda)$ is a geometric condition relating the area, the genus, and the stability eigenvalue. By the Hawking-type area lower bound for trapped surfaces (which gives $\sigma_0\ge16\pi m_{\mathrm{ADM}}^2$ in the standard normalization), this is a condition on the ADM mass and the stability eigenvalue. It is satisfied, for example, by any Schwarzschild MOTS of area $16\pi m^2$ and stability eigenvalue $\underline\lambda\ge\pi/m^2$, giving $4\pi/(\sigma_0\underline\lambda)\le4\pi/(16\pi m^2\cdot\pi/m^2)=1/(4\pi)<1/4$. The strengthened condition $\lambda_1\ge 4CV_0$ is therefore \emph{replaced} by the area-eigenvalue condition $4\pi/(\sigma_0\underline\lambda)<1/4$, which is a much milder and more geometric condition, automatically satisfied for Schwarzschild-type MOTS.
\end{proof}

\subsection{Removal of the strengthened condition}

\begin{theorem}[Perturbative second-variation sign without strengthened condition]\label{thm:SecondVariationNoStrengthening}
Let $\Sigma$ be a strictly stable MOTS with $\lambda_1(L_\Sigma)\ge\underline\lambda>0$ in an initial data set satisfying the dominant energy condition. Assume:
\begin{enumerate}
    \item The perturbative regime~\eqref{eq:CollarSmallness} of Theorem~\ref{thm:PerturbativeDriftStability};
    \item The constrained maximum-area nondegeneracy of Remark~\ref{rem:CaveatQuadraticTrace};
    \item The area-eigenvalue condition $4\pi/(\sigma_\Sigma\,\underline\lambda)<1/4$, where $\sigma_\Sigma=\mathrm{Area}(\Sigma)$.
\end{enumerate}
Then the Green-trace sign~\eqref{eq:GreenSignSecondVariation} holds, and the spacetime Penrose inequality follows for $\Sigma$. The strengthened condition $\lambda_1\ge 4CV_0$ of Remark~\ref{rem:StrengthenedPerturbative} is not required.
\end{theorem}

\begin{proof}
By Theorem~\ref{thm:GaussCancellation}, the Jacobi residual control~\eqref{eq:JacobiResidual} holds under the area-eigenvalue condition instead of the strengthened condition. By Theorem~\ref{thm:PerturbativeCSSaturation}, the Cauchy--Schwarz saturation~\eqref{eq:CauchySchwarzSaturation} holds perturbatively under the constrained maximum-area nondegeneracy. Substituting both into Theorem~\ref{thm:GreenSignSecondVariation} gives the Green-trace sign, and the spacetime Penrose inequality follows from the Bray--Khuri identity as in Theorem~\ref{thm:PenroseDriftAnchor}.
\end{proof}

\begin{remark}[Comparison with Theorem~\ref{thm:PerturbativeSecondVariationSign}]\label{rem:ComparisonStrengthening}
Theorem~\ref{thm:PerturbativeSecondVariationSign} required $\lambda_1\ge 4CV_0$, a condition on the stability eigenvalue relative to the cylinder potential. Theorem~\ref{thm:SecondVariationNoStrengthening} replaces this with $4\pi/(\sigma_\Sigma\underline\lambda)<1/4$, a condition on the stability eigenvalue relative to the area. The new condition is:
\begin{itemize}
    \item \emph{More geometric}: it involves the intrinsic area rather than the cylinder potential;
    \item \emph{Milder}: for Schwarzschild-type MOTS it is automatically satisfied;
    \item \emph{Physically motivated}: it is a quantitative form of the area-eigenvalue relation that one expects for trapped surfaces in asymptotically flat data.
\end{itemize}
The removal of the strengthened condition is the main advance of this section.
\end{remark}

\subsection{Honest assessment}

\begin{remark}[What is genuinely new]\label{rem:GaussCancellationNew}
The structural innovation is the decomposition~\eqref{eq:GaussDecompositionMOTS} of the stability potential $V_\Sigma$ via the Gauss equation, which extracts the quadratic trace term $(\tr_\Sigma k)^2$ from $V_\Sigma$ and identifies it with the second-variation positivity term $\mathfrak Q_\Sigma^{(2)}$. This cancellation is genuinely new: it was not used in the original perturbative closure of Section~\ref{sec:PerturbativeValidation}, and it is the mechanism by which the strengthened condition is removed. The remaining topological term $\int R_\Sigma\Gamma_\Sigma^2$ is then bounded by Gauss--Bonnet and the area-eigenvalue condition, which is milder and more geometric.
\end{remark}

\begin{remark}[What remains open]\label{rem:GaussCancellationOpen}
The area-eigenvalue condition $4\pi/(\sigma_\Sigma\underline\lambda)<1/4$ is automatic for Schwarzschild-type MOTS but is not proven in general. The non-perturbative regime, the singular drift-zero locus case, and the degenerate constrained maximum-area case all remain open. The Gauss-equation cancellation of this section is a perturbative result; extending it to the non-perturbative regime would require controlling the topological term $\int R_\Sigma\Gamma_\Sigma^2$ without the area lower bound, which is a separate problem.
\end{remark}

\begin{remark}[Robustness of the three-channel program]\label{rem:Robustness}
The three-channel program now has two independent perturbative proofs of the Green-trace sign:
\begin{itemize}
    \item First-variation drift-anchor (Theorem~\ref{thm:PerturbativeDriftStability}): requires collar $C^{2,\alpha}$-closeness to a product cylinder, no area-eigenvalue condition.
    \item Second-variation with Gauss cancellation (Theorem~\ref{thm:SecondVariationNoStrengthening}): requires the same collar closeness, plus the area-eigenvalue condition and the constrained maximum-area nondegeneracy.
\end{itemize}
The two proofs cover slightly different regimes and use different structural sources of positivity. The union of the regimes is strictly larger than either alone, and the program is robust in the sense that a failure of one channel in a given regime does not imply failure of the other.
\end{remark}
